% STATUS: written 2026-11-05 (sources of the section read in full: Ikeda & Kawamoto 1990 pp. 47--55; Wollenberg 1969 pp. 315--320; Li 2019 pp. 3250--3262; Randriambololondrantomalala--Asif 2022 pp. 51--65)
% STATUS: extended 2026-11-07 (Iena, Petravchuk and Regeta 2009 and
% Bezushchak and Kashuba 2025 read in full)
% Sources: IkedaKawamoto1990, Wollenberg1969, Li2019, RandriambololondrantomalalaAsif2022,
% BishtRaniTantubay2026 (Theorem 17--18, p. 18 and Corollary 24, p. 23, read 2026-11-06);
% Iena, Petravchuk and Regeta 2009 (pp. 1--4); Bezushchak and Kashuba 2025
% (pp. 1--2).
\section{Lie algebras of vector fields}\label{sec:lie-algebras:vector-fields}

Section~\ref{sec:lie-algebras:simple} showed that simple modular Lie
algebras can grow outer derivations.  Over fields of characteristic zero
some infinite-dimensional simple Lie algebras do the same
\cite{Wollenberg1969,IkedaKawamoto1990,Li2019}.  The examples
here are the polynomials under the Poisson bracket and the algebras of
vector fields studied in
\cite{Wollenberg1969,IkedaKawamoto1990,Li2019,RandriambololondrantomalalaAsif2022}.

\begin{definition}[{\cite[p.~315]{Wollenberg1969}}]\label{def:lie:poisson-bracket}
Let $E$ denote the collection of all $C^{\infty}$ complex-valued functions
of $2n$ real variables $x=(q_{1},\dots,q_{n},p_{1},\dots,p_{n})$, and
define the \emph{Poisson bracket} of two elements $f,g\in E$ by
\[
\{f,g\}=\sum_{j=1}^{n}\Big(\frac{\partial f}{\partial q_{j}}\frac{\partial g}{\partial p_{j}}-\frac{\partial f}{\partial p_{j}}\frac{\partial g}{\partial q_{j}}\Big).
\]
The collection of all complex polynomials in the variables $x$ forms a Lie
subalgebra of $E$ under this bracket, denoted $P$.
\end{definition}

The first observation of Wollenberg is that some derivations of $P$ are
not inner \cite[Lemma~1, p.~316]{Wollenberg1969}.

\begin{lemma}[{\cite[Lemma~1, p.~316]{Wollenberg1969}}]\label{lem:lie:wollenberg-outer}
For any set of complex numbers $\alpha=(\alpha_{1},\dots,\alpha_{n})$, the
map
\[
D_{\alpha}\colon f\longmapsto f\;-\;\sum_{i=1}^{n}\Big(\alpha_{i}p_{i}\frac{\partial f}{\partial p_{i}}+(1-\alpha_{i})q_{i}\frac{\partial f}{\partial q_{i}}\Big)
\]
is an outer derivation of $P$.
\end{lemma}

Together with the inner ones these give all derivations of $P$
\cite[Theorem~1, p.~318]{Wollenberg1969}.

\begin{theorem}[{\cite[Theorem~1, p.~318]{Wollenberg1969}}]\label{thm:lie:wollenberg-derivations}
Every derivation $D$ of $P$ is a linear combination of a derivation of
the form $D_{\alpha}$ of Lemma~\ref{lem:lie:wollenberg-outer} and an inner
derivation.  That is, one can find a polynomial $A$ and complex numbers
$c,\alpha_{i}$ such that $D(f)=c\,D_{\alpha}(f)+\{A,f\}$ for all $f\in P$.
\end{theorem}

The decomposition is not unique \cite[Remark, p.~318]{Wollenberg1969}:
shifting the $\alpha_{i}$ and changing $A$ give the same derivation.
Wollenberg also pins down the inner ones (Lemma~\ref{lem:lie:wollenberg-inner})
\cite[Lemma~3, p.~317]{Wollenberg1969}.\begin{lemma}[{\cite[Lemma~3, p.~317]{Wollenberg1969}}]\label{lem:lie:wollenberg-inner}
If $D$ is a derivation of $P$ with $D(1)=0$, then $D$ is inner.
\end{lemma}

A second family is the generalized Witt family of Ikeda and Kawamoto
\cite[p.~47]{IkedaKawamoto1990}.  Its index set $G$ is closed under
addition.

\begin{definition}[{\cite[p.~47]{IkedaKawamoto1990}}]\label{def:lie:generalized-witt}
Let $f$ be a field of characteristic zero, $I$ a non-empty index set, and
$G$ an additive submonoid of $\bigoplus_{i\in I} f t_{i}$ with
$t_{i}$ copies of the additive group $f^{+}$.  The \emph{generalized Witt
algebra} $W(G,I)$ is the Lie algebra over $f$ with basis
$\{w(a,i)\mid a\in G, i\in I\}$ and bracket
$[w(a,i),w(b,j)]=a_{j}w(a+b,i)-b_{i}w(a+b,j)$.
\end{definition}

For this algebra Ikeda and Kawamoto determine the full derivation
algebra (Theorem~\ref{thm:lie:generalized-witt-inner}) \cite[Theorem~2, p.~54]{IkedaKawamoto1990}.

\begin{theorem}[{\cite[Theorem~2, p.~54]{IkedaKawamoto1990}}]\label{thm:lie:generalized-witt-inner}
Let $G=\bigoplus_{i\in I}\mathbb{Z}_{i}$ and let $W=W(G,I)$.  Then $W$ is
simple and $\operatorname{Der} W=\operatorname{Lin} W$.  In particular, if
$|I|$ is finite and $G=\mathbb{Z}^{I}$, then
$\operatorname{Der} W=\operatorname{Inn} W$.
\end{theorem}

Here $\operatorname{Lin} W$ denotes the \emph{locally inner}
derivations of $W$ \cite[p.~47]{IkedaKawamoto1990}:
a derivation $\delta$ is locally inner when every finite subset
$F\subset W$ lies in a finite-dimensional subspace $V$ with
$\delta|_{V}=\operatorname{ad} x|_{V}$ \cite[p.~47]{IkedaKawamoto1990}.
In the finite-rank case every locally
inner derivation is inner; in infinite rank this fails in general
\cite[Remark, p.~55]{IkedaKawamoto1990}, giving the first outer
derivations of the section.

Li settles a parallel case \cite[Theorem~3.6]{Li2019}: the ring
$\mathbb{C}[z_{1}^{\pm 1},\dots,z_{n}^{\pm 1}]$ of Laurent polynomials,
its divergence-zero derivations, and once more an answer close to inner.

\begin{theorem}[{\cite[Theorem~3.6]{Li2019}}]\label{thm:lie:li-derivations}
Let $A=\mathbb{C}[z_{1}^{\pm 1},\dots,z_{n}^{\pm 1}]$ and let $S_{n}$ be
the Lie algebra of divergence-zero derivations of $A$.  Let
$L=[S_{n}\!\cdot\!A,S_{n}\!\cdot\!A]$ be the derived subalgebra of the
semi-direct product.  Then
$\operatorname{Der} L=\operatorname{Inn} L\oplus\operatorname{Der}(L)_{0}$,
where $\operatorname{Der}(L)_{0}$ has dimension $2n+1$ and is spanned by
explicit degree derivations.
\end{theorem}

In this decomposition the extra summand $\operatorname{Der}(L)_{0}$
counts the degree derivations and one extra direction
\cite[Theorem~3.6]{Li2019}.  On the remaining weight spaces every
derivation is inner \cite[Theorem~3.6]{Li2019}.  Once more the outer
derivations come from the grading data, as in the Poisson case of
Theorem~\ref{thm:lie:wollenberg-derivations}.
Randriambololondrantomalala and Asif give an explicit description of
this quotient $\operatorname{Der}(P)/\operatorname{ad}_{P}$
\cite[Theorem~2.4, p.~52]{RandriambololondrantomalalaAsif2022}.

\begin{theorem}[{\cite[Theorem~2.4, p.~52]{RandriambololondrantomalalaAsif2022}}]\label{thm:lie:h1-vector-fields}
Let $P$ be a Lie subalgebra of the Lie algebra $\chi(\mathbb{R}^{n})$ of
vector fields on $\mathbb{R}^{n}$, containing all constant vector fields
and the Euler vector field, and write
$H^{1}(P):=\operatorname{Der}(P)/\operatorname{ad}_{P}$.  Then
$H^{1}(P)$ is isomorphic to $N/P$, where $N$ is the normaliser of $P$ in
$\chi(\mathbb{R}^{n})$; if all linear diagonal fields lie in $P$, then
$H^{1}(P)=0$.
\end{theorem}

Together Theorem~\ref{thm:lie:h1-vector-fields} identifies all outer
derivations of such a $P$ with an explicit quotient, and
Theorems~\ref{thm:lie:wollenberg-derivations}--\ref{thm:lie:li-derivations}
compute it in the classical examples.

The Hamiltonian Lie algebras on the torus give a case where the answer is
inner. Let $N$ be an even positive integer, and let $H_{N}$ be the
Hamiltonian Lie algebra on the torus, with derived subalgebra $H_{N}'$
\cite[p.~1]{BishtRaniTantubay2026}. Bisht, Rani and Tantubay compute both
derivation algebras (Theorem~\ref{thm:lie:hamiltonian-inner}) \cite[Theorems~17--18, p.~18]{BishtRaniTantubay2026}.

\begin{theorem}[{\cite[Theorems~17--18, p.~18]{BishtRaniTantubay2026}}]\label{thm:lie:hamiltonian-inner}
Let $N$ be an even positive integer. Then
$\operatorname{Der}(H_{N}')=H_{N}'\rtimes\mathfrak{h}\cong H_{N}$ and
$\operatorname{Der}(H_{N})\cong H_{N}$.
\end{theorem}

So every derivation of the Hamiltonian Lie algebra is inner, and the derived
subalgebra has one more family, indexed by the diagonal fields $\mathfrak{h}$
\cite[p.~18]{BishtRaniTantubay2026}. The authors also compute the second
cohomology of $H_{N}$ \cite[Corollary~24, p.~23]{BishtRaniTantubay2026}.

The same question has a geometric answer on the plane. Let $k$ be an
algebraically closed field of characteristic zero and let $W_2$ be the Lie
algebra of all $k$-derivations of $k[x,y]$ \cite[p.~1]{IenaPetravchukRegeta2009}.
For a fixed rational function $u\in k(x,y)\setminus k$, the annihilator
$A_{W_2}(u)$ is the set of derivations $D\in W_2$ with $D(u)=0$
\cite[p.~2]{IenaPetravchukRegeta2009}. It is a Lie subalgebra of $W_2$ and a
$k[x,y]$-submodule of the module $W_2$ (Theorem~\ref{thm:lie:annihilator-free}) \cite[p.~2]{IenaPetravchukRegeta2009}.

\begin{theorem}[{\cite[Theorem~6, p.~4]{IenaPetravchukRegeta2009}}]\label{thm:lie:annihilator-free}
For an arbitrary rational function $\phi\in k(x,y)\setminus k$, its annihilator
$A_{W_2}(\phi)$ is a free submodule of rank one of the $k[x,y]$-module $W_2(k)$.
As a free generator one may take the derivation $\delta_\phi$.
\end{theorem}

The generator $\delta_\phi$ is reduced, meaning its two polynomial coefficients
are coprime \cite[p.~2]{IenaPetravchukRegeta2009}. The paper also describes the
centralisers of elements of $A_{W_2}(\phi)$ and its maximal abelian
subalgebras \cite[p.~2]{IenaPetravchukRegeta2009}. The annihilator itself has a
different structure for a polynomial and for a quotient of two algebraically
independent polynomials \cite[p.~2]{IenaPetravchukRegeta2009}.

The polynomial ring in infinitely many variables gives an algebra whose
derivations are all inner. Let $F$ be a field of characteristic zero and let
$W_X(F)$ be the Lie algebra of all derivations of $F[X]$ in infinitely many
variables \cite[p.~1]{BezushchakKashuba2025}. Two subalgebras sit inside it,
$W_\infty(F)$ and $W_{\mathrm{fin}}(F)$ (Theorem~\ref{thm:lie:infinitely-many-inner}) \cite[p.~2]{BezushchakKashuba2025}.

\begin{theorem}[{\cite[Theorems~1--3, p.~2]{BezushchakKashuba2025}}]\label{thm:lie:infinitely-many-inner}
Let $F$ be a field of characteristic zero. Every derivation of $W_\infty(F)$ is
the restriction of an adjoint operator $\operatorname{ad}(a)$ with
$a\in W_{\mathrm{fin}}(F)$. Every derivation of $W_{\mathrm{fin}}(F)$ is inner,
and every derivation of $W_X(F)$ is inner.
\end{theorem}

So $W_\infty(F)$ has derivations that are not inner, and its inner derivations
are dense in its derivations for the Tychonoff topology
\cite[p.~2]{BezushchakKashuba2025}. The algebra $W_{\mathrm{fin}}(F)$ and the
full algebra $W_X(F)$ have no outer derivation at all
\cite[p.~2]{BezushchakKashuba2025}. This differs from the Poisson case of
Theorem~\ref{thm:lie:wollenberg-derivations}, where the degree derivation gave
an outer one \cite[Lemma~1, p.~316]{Wollenberg1969}.

Section~\ref{sec:lie-algebras:infinite-dimensional}
pushes the same question into Kac--Moody theory, where the algebra is no
longer a vector-field algebra but an affine extension of a finite
semisimple one.

Further work in the Lie-algebra line, held in the collection:
\noindent\emph{The groups of automorphisms of the Lie algebras of polynomial vector fields with zero or constant divergence}~\cite{LibTheGroupsOfAutomorphisms2017}.
\noindent\emph{Space of second-order linear differential operators as a module over the Lie algebra of vector fields}~\cite{LibSpaceOfSecondOrder1997}.
\noindent\emph{Classification of finite-dimensional Lie superalgebras whose even part is a three-dimensional simple Lie algebra over a field of characteristic not two or three}~\cite{LibClassificationOfFiniteDimensional2019}.
\noindent\emph{The groups of automorphisms of the Lie algebras of formally analytic vector fields with constant divergence}~\cite{LibTheGroupsOfAutomorphisms2014}.

On the transposed Poisson algebra $W(a,-1)$ the derivations, local derivations, quasi-derivations and $\delta$-derivations are all classified, and the nontrivial ones are separated from the trivial~\cite{WlTransposedPoissonStructure2026}.

For a smooth foliated manifold the Lie algebra of leaf-tangent vector fields is studied for the partially classical case, and its derivations are determined~\cite{WlSomeLieAlgebras1981}.

The nongraded Hamiltonian Lie algebras that Xu introduced carry a Poisson bracket structure, and a sandwich argument determines their isomorphism classes~\cite{WlPoissonBracketsAnd2003}.

Basdouri, Hamdouni, Nasri and Missaoui compute the first cohomology space of the Lie algebra of vector fields on the affine real line, relative to the affine vector fields, with coefficients in bilinear differential operators on weighted densities~\cite{WlFirstCohomologyOfTheLieAlgebra2020}.

Ficzere and Figula determine the systems of first order ordinary differential equations whose symmetry group contains a given three-dimensional Lie group~\cite{WlSystemsOfFirstOrderOrdinaryDif2020}.

Billig, Molev and Zhang describe differential equations in vertex algebras and use them to construct simple modules for the Lie algebra of vector fields on a torus~\cite{WlBilligYMolevAZhangRDifferentia2020}.

Post describes a class of graded Lie algebras of vector fields together with their first order differential operators~\cite{WlAClassOfGradedLieAlgebrasOfVec2020}.

Torossian studies invariant differential operators on symmetric spaces through a generalised Harish-Chandra homomorphism, the continuation of the paper on the same subject~\cite{WlInvariantDifferentialOperators2020}.

The cohomology of the Lie algebras of vector fields on the real line acting on trilinear differential operators is computed~\cite{WlCohomologyOfTheVectorFieldsLie2020}.

Bouarroudj computes the cohomology of the Lie algebras of vector fields with coefficients in modules of differential operators~\cite{WlBouarroudjSofianeCohomologyOfT2020}.

The cohomology of the Lie algebras of vector fields on the real line acting on trilinear differential operators is computed~\cite{WlCohomologyOfTheVectorFieldsLie2020}.

The cohomology of the Lie algebras of vector fields on the projective line acting on differential operators is computed~\cite{WlCohomologyOfTheVectorFieldsLie2020}.

Bouarroudj computes the relative cohomology of the Lie algebra $\mathfrak{sl}(2)$ with coefficients in differential operator modules~\cite{WlBouarroudjSOnMathfrakSl2Relati2020}.

Invariant differential operators on the tangent space of some symmetric spaces are described~\cite{WlInvariantDifferentialOperators2020}.
